\documentclass[12pt,a4paper]{article}
\usepackage[utf8]{inputenc}
\usepackage{amsmath,amsfonts,amssymb}
\usepackage{listings}
\usepackage{xcolor}
\usepackage{geometry}
\usepackage{hyperref}
\usepackage{fancyhdr}

\geometry{margin=1in}
\pagestyle{fancy}
\fancyhf{}
\rhead{Python Sets - Complete Guide}
\lfoot{Python for Data Science}
\rfoot{\thepage}

\lstset{
    language=Python,
    basicstyle=\ttfamily\small,
    keywordstyle=\color{blue}\bfseries,
    commentstyle=\color{green!60!black},
    stringstyle=\color{red},
    numbers=left,
    numberstyle=\tiny\color{gray},
    stepnumber=1,
    numbersep=8pt,
    showstringspaces=false,
    breaklines=true,
    frame=single,
    backgroundcolor=\color{gray!10}
}

\title{\textbf{Python Sets - Complete Guide}}
\author{Python for Data Science - Beginner's Level}
\date{\today}

\begin{document}

\maketitle
\tableofcontents
\newpage

\section{Overview}

Sets are \textbf{unordered}, \textbf{mutable} collections of \textbf{unique} elements. They are ideal for membership testing, removing duplicates, and mathematical set operations like union, intersection, and difference.

\section{Set Creation}

\subsection{Basic Creation}

\begin{lstlisting}
# Empty set (note: {} creates an empty dict, not set)
empty_set = set()

# Set with initial values
numbers = {1, 2, 3, 4, 5}
mixed_set = {1, "hello", 3.14, True}

# Note: duplicates are automatically removed
duplicate_values = {1, 2, 2, 3, 3, 4}
print(duplicate_values)  # {1, 2, 3, 4}

# Set from string (each character becomes an element)
char_set = set("hello")
print(char_set)  # {'h', 'e', 'l', 'o'}
\end{lstlisting}

\subsection{Using set() Constructor}

\begin{lstlisting}
# From list
from_list = set([1, 2, 3, 2, 1])  # {1, 2, 3}

# From tuple
from_tuple = set((1, 2, 3, 2, 1))  # {1, 2, 3}

# From string
from_string = set("programming")  # {'p', 'r', 'o', 'g', 'a', 'm', 'i', 'n'}

# From range
from_range = set(range(5))  # {0, 1, 2, 3, 4}
\end{lstlisting}

\subsection{Set Comprehension}

\begin{lstlisting}
# Basic comprehension
squares = {x**2 for x in range(6)}  # {0, 1, 4, 9, 16, 25}

# With condition
even_squares = {x**2 for x in range(10) if x % 2 == 0}

# From existing iterable
words = ["apple", "banana", "apple", "cherry"]
unique_lengths = {len(word) for word in words}  # {5, 6}
\end{lstlisting}

\section{Adding and Removing Elements}

\subsection{Adding Elements}

\begin{lstlisting}
fruits = {"apple", "banana"}

# add() - Add single element
fruits.add("cherry")

# Adding existing element (no effect)
fruits.add("apple")  # No change

# update() - Add multiple elements
fruits.update(["orange", "grape"])

# Using |= operator (union update)
more_fruits = {"pineapple", "strawberry"}
fruits |= more_fruits
\end{lstlisting}

\subsection{Removing Elements}

\begin{lstlisting}
fruits = {"apple", "banana", "cherry", "orange", "grape"}

# remove() - Remove element (raises KeyError if not found)
fruits.remove("banana")

# discard() - Remove element (no error if not found)
fruits.discard("cherry")
fruits.discard("mango")  # No error

# pop() - Remove and return arbitrary element
removed = fruits.pop()

# clear() - Remove all elements
fruits.clear()
\end{lstlisting}

\section{Set Operations and Methods}

\subsection{Membership Testing}

\begin{lstlisting}
numbers = {1, 2, 3, 4, 5}

# in operator (very fast - O(1) average case)
print(3 in numbers)      # True
print(6 in numbers)      # False
print(3 not in numbers)  # False
\end{lstlisting}

\subsection{Mathematical Set Operations}

\subsubsection{Union Operations}

\begin{lstlisting}
set1 = {1, 2, 3, 4}
set2 = {3, 4, 5, 6}

# union() method
union_result = set1.union(set2)  # {1, 2, 3, 4, 5, 6}

# | operator
union_operator = set1 | set2  # {1, 2, 3, 4, 5, 6}

# In-place union
set1_copy = set1.copy()
set1_copy.update(set2)  # or set1_copy |= set2
\end{lstlisting}

\subsubsection{Intersection Operations}

\begin{lstlisting}
set1 = {1, 2, 3, 4, 5}
set2 = {3, 4, 5, 6, 7}

# intersection() method
intersection_result = set1.intersection(set2)  # {3, 4, 5}

# & operator
intersection_operator = set1 & set2  # {3, 4, 5}

# In-place intersection
set1_copy = set1.copy()
set1_copy.intersection_update(set2)  # or set1_copy &= set2
\end{lstlisting}

\subsubsection{Difference Operations}

\begin{lstlisting}
set1 = {1, 2, 3, 4, 5}
set2 = {3, 4, 5, 6, 7}

# difference() method (elements in set1 but not in set2)
difference_result = set1.difference(set2)  # {1, 2}

# - operator
difference_operator = set1 - set2  # {1, 2}

# Reverse difference
reverse_diff = set2 - set1  # {6, 7}
\end{lstlisting}

\subsubsection{Symmetric Difference Operations}

\begin{lstlisting}
set1 = {1, 2, 3, 4}
set2 = {3, 4, 5, 6}

# symmetric_difference() method (elements in either set, but not both)
sym_diff_result = set1.symmetric_difference(set2)  # {1, 2, 5, 6}

# ^ operator
sym_diff_operator = set1 ^ set2  # {1, 2, 5, 6}
\end{lstlisting}

\section{Set Relationship Testing}

\begin{lstlisting}
set1 = {1, 2, 3}
set2 = {1, 2, 3, 4, 5}
set3 = {4, 5, 6}

# issubset() - Check if set is subset of another
print(set1.issubset(set2))  # True
print(set1 <= set2)         # True (alternative syntax)

# issuperset() - Check if set is superset of another
print(set2.issuperset(set1))  # True
print(set2 >= set1)           # True (alternative syntax)

# isdisjoint() - Check if sets have no common elements
print(set1.isdisjoint(set3))  # True
\end{lstlisting}

\section{Advanced Set Operations}

\subsection{Frozen Sets}

\begin{lstlisting}
# Immutable sets - can be used as dictionary keys
regular_set = {1, 2, 3}
frozen_set = frozenset([1, 2, 3])

# Frozen sets can be dictionary keys
dict_with_frozenset_keys = {
    frozenset([1, 2]): "first",
    frozenset([3, 4]): "second"
}

# Frozen sets support all non-mutating operations
fs1 = frozenset([1, 2, 3])
fs2 = frozenset([2, 3, 4])

print(fs1 | fs2)  # frozenset({1, 2, 3, 4})
print(fs1 & fs2)  # frozenset({2, 3})
\end{lstlisting}

\section{Performance Considerations}

\begin{table}[h]
\centering
\begin{tabular}{|l|c|c|l|}
\hline
\textbf{Operation} & \textbf{Average} & \textbf{Worst} & \textbf{Notes} \\
\hline
Add & O(1) & O(n) & Hash collision \\
Remove & O(1) & O(n) & Hash collision \\
Search (in) & O(1) & O(n) & Hash collision \\
Union & O(len(s1) + len(s2)) & O(len(s1) * len(s2)) & \\
Intersection & O(min(len(s1), len(s2))) & O(len(s1) * len(s2)) & \\
\hline
\end{tabular}
\caption{Set Operations Time Complexity}
\end{table}

\section{Set Applications and Use Cases}

\subsection{Removing Duplicates}

\begin{lstlisting}
# Remove duplicates from list
numbers = [1, 2, 2, 3, 1, 4, 3, 5]
unique_numbers = list(set(numbers))

# Preserve order while removing duplicates
def remove_duplicates_ordered(lst):
    seen = set()
    result = []
    for item in lst:
        if item not in seen:
            seen.add(item)
            result.append(item)
    return result
\end{lstlisting}

\subsection{Finding Common Elements}

\begin{lstlisting}
# Find common elements across multiple lists
def find_common_elements(*lists):
    if not lists:
        return set()
    
    common = set(lists[0])
    for lst in lists[1:]:
        common &= set(lst)
    return common

# Test
list1 = [1, 2, 3, 4, 5]
list2 = [3, 4, 5, 6, 7]
list3 = [4, 5, 6, 7, 8]

common = find_common_elements(list1, list2, list3)  # {4, 5}
\end{lstlisting}

\subsection{Data Validation}

\begin{lstlisting}
# Validate data against allowed values
ALLOWED_COLORS = {"red", "green", "blue", "yellow", "orange"}
ALLOWED_SIZES = {"small", "medium", "large", "xl", "xxl"}

def validate_product(color, size):
    errors = []
    
    if color.lower() not in ALLOWED_COLORS:
        errors.append(f"Invalid color: {color}")
    
    if size.lower() not in ALLOWED_SIZES:
        errors.append(f"Invalid size: {size}")
    
    return errors
\end{lstlisting}

\section{Common Pitfalls}

\subsection{Mutable Elements in Sets}

\begin{lstlisting}
# This won't work - lists are mutable and unhashable
# invalid_set = {[1, 2], [3, 4]}  # TypeError

# Use tuples or frozensets instead
valid_set = {(1, 2), (3, 4)}
frozenset_set = {frozenset([1, 2]), frozenset([3, 4])}
\end{lstlisting}

\subsection{Empty Set Creation}

\begin{lstlisting}
# Wrong way to create empty set
empty_dict = {}  # This creates a dictionary, not a set!

# Right way to create empty set
empty_set = set()
\end{lstlisting}

\section{Summary}

Sets are powerful Python data structures that provide:
\begin{itemize}
    \item \textbf{Uniqueness}: Automatically handle duplicate elimination
    \item \textbf{Fast membership testing}: O(1) average case lookup
    \item \textbf{Mathematical operations}: Union, intersection, difference, etc.
    \item \textbf{Unordered collections}: Focus on membership, not position
    \item \textbf{Hashable elements only}: Immutable objects as elements
    \item \textbf{Memory efficiency}: Optimized for unique element storage
\end{itemize}

Master sets to efficiently handle unique collections and perform mathematical set operations!

\end{document}